% Сообщение ученику: Ярослав, в работе остались два момента: в задании 1в решение не было доведено до записи через косинус половинного аргумента, а в задании 3 после верного получения косинуса икс, равного одной второй, был выбран неверный угол: вместо пи третьих записано пи шестых. Вторая ветвь решения уравнения, где косинус икс равен минус единице, записана верно.
\documentclass[12pt,a4paper]{article}
\usepackage[a4paper,margin=2cm]{geometry}
\usepackage[T2A]{fontenc}
\usepackage[utf8]{inputenc}
\usepackage[russian]{babel}
\usepackage{amsmath,amssymb,mathtools}
\usepackage{array,booktabs,longtable}
\usepackage{xcolor}
\usepackage{hyperref}
\usepackage{tikz}

\hypersetup{
  colorlinks=true,
  linkcolor=blue!55!black,
  urlcolor=blue!55!black,
  pdftitle={Ревью домашней работы},
  pdfauthor={Лёвин Артём Александрович}
}

\setlength{\parindent}{0pt}
\setlength{\parskip}{0.65em}
\renewcommand{\arraystretch}{1.25}

\newcommand{\errorlabel}[1]{\textcolor{red}{Ошибка:} #1}
\newcommand{\keyidea}[1]{\textcolor{blue!60!black}{Ключевая идея:} #1}
\newcommand{\resultlabel}[1]{\textcolor{teal!60!black}{Итог:} #1}

\begin{document}

\begin{titlepage}
\centering
\vspace*{0.22\textheight}
{\LARGE\bfseries Ревью домашней работы\par}
\vspace{2cm}
{\large Автор: Лёвин Артём Александрович\par}
\vspace{0.8cm}
{\large 24 сентября 2026 года\par}
\vfill
\begin{tikzpicture}
\draw[blue!35!black, line width=0.9pt]
  (0,0) .. controls (2.5,0.55) and (5,-0.55) .. (7.5,0);
\end{tikzpicture}
\end{titlepage}

\newpage
\section*{Сводная таблица}
\small
\setlength{\tabcolsep}{3pt}
\begin{longtable}{>{\centering\arraybackslash}p{0.08\linewidth} p{0.17\linewidth} p{0.27\linewidth} p{0.18\linewidth} p{0.18\linewidth}}
\toprule
№ задания & Тема & Суть ошибки & Ваш ответ & Правильный ответ \\
\midrule
\endfirsthead
\toprule
№ задания & Тема & Суть ошибки & Ваш ответ & Правильный ответ \\
\midrule
\endhead
\hyperref[task:1c]{1в} & Формула двойного угла & Решение не выполнено: в работе записано «не понял», поэтому требуемое представление не получено. & не указан & $2\cos^2\!\left(\frac{3x}{2}\right)-1$ \\
\addlinespace
\hyperref[task:3]{3} & Тригоно\-метри\-ческое уравнение; сведение к квадратному & После верного получения $\cos x=\frac12$ выбран неверный угол: записано $\frac{\pi}{6}$ вместо $\frac{\pi}{3}$. & $x=\pm\frac{\pi}{6}+2\pi k$;\newline $x=\pi+2\pi k$, $k\in\mathbb Z$ & $x=\pm\frac{\pi}{3}+2\pi k$;\newline $x=\pi+2\pi k$, $k\in\mathbb Z$ \\
\bottomrule
\end{longtable}
\normalsize

\newpage
\thispagestyle{empty}
\vspace*{\fill}
\begin{center}
{\LARGE\bfseries Разбор ошибок}
\end{center}
\vspace*{\fill}

\newpage
\phantomsection\label{task:1c}
\section*{Задание 1в}

\subsection*{Текст задания}
Представьте $\cos 3x$ только через $\cos\!\left(\frac{3x}{2}\right)$.

\subsection*{Ваше решение}
В работе записано: «не понял». Математические преобразования отсутствуют.

\subsection*{Ваш ответ}
Не указан.

\subsection*{Ошибка}
\errorlabel{решение не выполнено.}

\subsection*{Где именно возникает ошибка}
Последнего корректного математического шага в записи нет: после условия не выполнено ни одного преобразования.

Первый неполный шаг состоит в том, что не замечено представление аргумента $3x$ как удвоенного аргумента:
\[
3x=2\cdot\frac{3x}{2}.
\]
Именно это позволяет применить формулу косинуса двойного угла.

\subsection*{Почему это неверно}
Требование «представьте только через $\cos\!\left(\frac{3x}{2}\right)$» означает, что в конечном выражении не должно оставаться $\cos 3x$ или других тригонометрических функций с иными аргументами. Нужно заменить косинус удвоенного аргумента выражением через косинус половинного аргумента.

Для любого угла $u$ справедлива школьная формула двойного угла
\[
\cos 2u=2\cos^2u-1.
\]
Здесь роль $u$ играет $\frac{3x}{2}$, потому что удвоение этого аргумента даёт $3x$.

\subsection*{Ключевая идея}
\keyidea{сначала представить $3x$ как $2\cdot\frac{3x}{2}$, затем применить формулу $\cos 2u=2\cos^2u-1$.}

\subsection*{Исправление от последнего правильного шага}
Так как математических шагов в исходной записи нет, исправление начинается непосредственно с выбора подходящей формулы:
\[
\cos 3x
=\cos\left(2\cdot\frac{3x}{2}\right).
\]
Теперь применяем формулу двойного угла:
\[
\cos\left(2\cdot\frac{3x}{2}\right)
=2\cos^2\!\left(\frac{3x}{2}\right)-1.
\]

\subsection*{Полное правильное решение}
Используем формулу
\[
\cos 2u=2\cos^2u-1.
\]
Положим
\[
u=\frac{3x}{2}.
\]
Тогда
\[
2u=3x.
\]
Следовательно,
\[
\cos 3x
=\cos 2u
=2\cos^2u-1
=2\cos^2\!\left(\frac{3x}{2}\right)-1.
\]

\subsection*{Проверка}
В полученном выражении используется только $\cos\!\left(\frac{3x}{2}\right)$ и числовые операции над ним. Обратная подстановка $u=\frac{3x}{2}$ в формулу двойного угла точно возвращает $\cos 3x$.

\subsection*{Итог}
\resultlabel{правильное представление:}
\[
\boxed{\cos 3x=2\cos^2\!\left(\frac{3x}{2}\right)-1}.
\]

\subsection*{Задача на закрепление}
Представьте $\cos 5x$ только через $\cos\!\left(\frac{5x}{2}\right)$.

\newpage
\phantomsection\label{task:3}
\section*{Задание 3}

\subsection*{Текст задания}
Решите уравнение
\[
\cos 2x+\sin\left(x+\frac{\pi}{2}\right)=0.
\]
Запишите все решения.

\subsection*{Ваше решение}
В работе верно выполнено преобразование
\[
\sin\left(x+\frac{\pi}{2}\right)=\cos x,
\]
после чего уравнение сведено к виду
\[
2\cos^2x+\cos x-1=0.
\]
Далее введена замена
\[
t=\cos x,
\]
и решено квадратное уравнение
\[
2t^2+t-1=0.
\]
Получены корни
\[
t_1=\frac12,\qquad t_2=-1.
\]
После этого для первого корня записано
\[
x=\pm\frac{\pi}{6}+2\pi k,
\]
а для второго
\[
x=\pi+2\pi k,\qquad k\in\mathbb Z.
\]

\subsection*{Ваш ответ}
\[
x=\pm\frac{\pi}{6}+2\pi k,
\qquad
x=\pi+2\pi k,
\qquad k\in\mathbb Z.
\]

\subsection*{Ошибка}
\errorlabel{неверно найден угол для уравнения $\cos x=\frac12$.}

\subsection*{Где именно возникает ошибка}
Последний корректный шаг:
\[
\cos x=\frac12.
\]

Первый неверный шаг:
\[
\cos x=\frac12
\quad\Longrightarrow\quad
x=\pm\frac{\pi}{6}+2\pi k.
\]

Этот переход неверен, потому что
\[
\cos\frac{\pi}{6}=\frac{\sqrt3}{2},
\]
а не $\frac12$.

\subsection*{Почему этот шаг неверен}
Для решения простейшего тригонометрического уравнения $\cos x=a$ нужно выбрать угол, косинус которого действительно равен числу $a$.

В данном случае
\[
\cos\frac{\pi}{3}=\frac12.
\]
Поэтому базовый угол равен $\frac{\pi}{3}$, а не $\frac{\pi}{6}$.

Косинус является чётной функцией, поэтому два семейства углов удобно объединить в запись
\[
x=\pm\frac{\pi}{3}+2\pi k,
\qquad k\in\mathbb Z.
\]

Ветка $\cos x=-1$ в работе решена правильно:
\[
x=\pi+2\pi k,
\qquad k\in\mathbb Z.
\]
Её изменять не нужно.

\subsection*{Ключевая идея}
\keyidea{после сведения уравнения к простейшим тригонометрическим уравнениям отдельно проверить табличное значение каждого найденного косинуса. Для $\cos x=\frac12$ опорный угол равен $\frac{\pi}{3}$.}

\subsection*{Исправление от последнего правильного шага}
Из уже правильно полученного уравнения
\[
\cos x=\frac12
\]
следует
\[
x=\pm\frac{\pi}{3}+2\pi k,
\qquad k\in\mathbb Z.
\]

Вторая ветвь сохраняется без изменений:
\[
\cos x=-1
\quad\Longrightarrow\quad
x=\pi+2\pi k,
\qquad k\in\mathbb Z.
\]

\subsection*{Полное правильное решение}
Начинаем с исходного уравнения:
\[
\cos 2x+\sin\left(x+\frac{\pi}{2}\right)=0.
\]

По формуле суммы для синуса
\[
\sin\left(x+\frac{\pi}{2}\right)
=\sin x\cos\frac{\pi}{2}+\cos x\sin\frac{\pi}{2}
=\cos x.
\]
Поэтому
\[
\cos 2x+\cos x=0.
\]

Используем формулу двойного угла
\[
\cos 2x=2\cos^2x-1.
\]
Получаем
\[
2\cos^2x-1+\cos x=0,
\]
то есть
\[
2\cos^2x+\cos x-1=0.
\]

Введём замену
\[
t=\cos x.
\]
Тогда
\[
2t^2+t-1=0.
\]
Дискриминант:
\[
D=1^2-4\cdot2\cdot(-1)=9.
\]
Следовательно,
\[
t_1=\frac{-1+3}{4}=\frac12,
\qquad
t_2=\frac{-1-3}{4}=-1.
\]
Возвращаемся к $\cos x$.

Первая ветвь:
\[
\cos x=\frac12
\quad\Longrightarrow\quad
x=\pm\frac{\pi}{3}+2\pi k,
\qquad k\in\mathbb Z.
\]

Вторая ветвь:
\[
\cos x=-1
\quad\Longrightarrow\quad
x=\pi+2\pi k,
\qquad k\in\mathbb Z.
\]

Итак, все решения:
\[
\boxed{x=\pm\frac{\pi}{3}+2\pi k
\quad\text{или}\quad
x=\pi+2\pi k,
\qquad k\in\mathbb Z.}
\]

\subsection*{Проверка}
Для
\[
x=\pm\frac{\pi}{3}+2\pi k
\]
имеем
\[
\cos x=\frac12,
\qquad
\cos 2x=-\frac12,
\qquad
\sin\left(x+\frac{\pi}{2}\right)=\cos x=\frac12.
\]
Поэтому левая часть исходного уравнения равна нулю.

Для
\[
x=\pi+2\pi k
\]
имеем
\[
\cos x=-1,
\qquad
\cos 2x=1,
\qquad
\sin\left(x+\frac{\pi}{2}\right)=\cos x=-1.
\]
Снова левая часть равна нулю.

\subsection*{Итог}
\resultlabel{ошибка возникла только при выборе табличного угла для значения $\cos x=\frac12$. Правильный угол --- $\frac{\pi}{3}$. Поэтому полный ответ:}
\[
\boxed{x=\pm\frac{\pi}{3}+2\pi k
\quad\text{или}\quad
x=\pi+2\pi k,
\qquad k\in\mathbb Z.}
\]

\subsection*{Задача на закрепление}
Решите уравнение
\[
2\cos^2x-3\cos x+1=0.
\]
Запишите все решения.

\vfill
\begin{center}
\small Лёвин Артём Александрович эксклюзивно для Ярослава Гаврилова
\end{center}

\end{document}
